% !TEX root = ../main.tex
% Resultados literales de la tanda definitiva, repuntuada en IT-110.
% Fuente: el bloque automatico de docs/adr/adr-0005-modelo-de-generacion.md,
% que escribe scripts/experimentos/experimento_generacion.py --adr.
\begin{table}[htbp]
  \centering
  \small
  \caption[Comparativa de modelos de generación]{%
    Comparativa de los modelos de generación medidos sobre
      la misma muestra de 80 preguntas, 320 respuestas en total. Fuente:
      elaboración propia.}
  \label{tab:comparativa_generacion}
  \begin{tabular}{@{}lrrrrrr@{}}
    \toprule
    \textbf{Modelo}                    & \textbf{Inv.} & \textbf{Precisión}
                                       & \textbf{Cobertura}
                                       & \textbf{Escalar}
                                       & \textbf{Mediana} & \textbf{p90}   \\
    \midrule
    \texttt{ministral-8b:latest}       & \textbf{0}    & \textbf{1,000}
                                       & \sout{0,996} \textcolor{borrador2}{0,971}
                                       & \sout{0,957} \textcolor{borrador2}{1,000}
                                       & \sout{12,1} \textcolor{borrador2}{24,4}
                                       & \sout{42,1} \textcolor{borrador2}{87,2}   \\
    \texttt{qwen3.5:9b}                & \textbf{0}    & \sout{0,996} \textcolor{borrador2}{\textbf{1,000}}
                                       & \sout{0,997} \textcolor{borrador2}{0,971}
                                       & \sout{0,957} \textcolor{borrador2}{1,000}
                                       & \sout{18,2} \textcolor{borrador2}{17,4}
                                       & \sout{67,2} \textcolor{borrador2}{47,7}   \\
    \texttt{gemma3:12b}                & \textbf{0}    & \textbf{1,000}
                                       & \sout{\textbf{1,000}} \textcolor{borrador2}{0,971}
                                       & \sout{\textbf{0,978}} \textcolor{borrador2}{1,000}
                                       & \sout{21,6} \textcolor{borrador2}{26,0}
                                       & \sout{74,4} \textcolor{borrador2}{72,1}   \\
    \midrule
    \texttt{salamandra-7b}             & \textbf{0}    & \sout{0,992} \textcolor{borrador2}{0,995}
                                       & \sout{0,908} \textcolor{borrador2}{0,770}
                                       & \sout{0,935} \textcolor{borrador2}{0,978}
                                       & \sout{3,8} \textcolor{borrador2}{7,1}
                                       & \sout{11,1} \textcolor{borrador2}{21,2}   \\
    \bottomrule
  \end{tabular}
  \begin{minipage}{0.95\textwidth}
    \vspace{4pt}
    \tiny
    \textbf{Notas:} \textbf{Inv.} son titulaciones inventadas (criterio eliminatorio).
    \textbf{Precisión} y \textbf{cobertura} se promedian sobre las 34 preguntas de listado.
    \textbf{Escalar} es el acierto sobre las 46 preguntas de créditos y curso.
    \textbf{Mediana} y \textbf{p90} indican la latencia en segundos (datos informativos).
    En negrita, el mejor valor por columna salvo en tiempos.
  \end{minipage}
\end{table}
